\documentclass[11pt,a4paper]{article}
\usepackage[utf8]{inputenc}
\usepackage[margin=1.0in]{geometry}
\usepackage{hyperref}
\usepackage{microtype}
\usepackage{charter}
\usepackage{xcolor}

\hypersetup{
    colorlinks=true,
    linkcolor=blue!70!black,
    urlcolor=blue!70!black
}

\title{\textbf{Ethics and Compliance Statement}}
\author{\textbf{Ronie Martínez}}
\date{November 6, 2026}

\begin{document}

\maketitle

\section*{Manuscript Details}
\textbf{Title:} Rigorous Benchmarking and De-Biasing of Machine Learning Predictors for Antimicrobial Peptides: Resolving Homology Leakage and Negative Sampling Artifacts\\
\textbf{Corresponding Author:} Ronie Martínez (\texttt{ronie.martinez@epn.edu.ec})\\
\textbf{Affiliation:} Departamento de Informática y Ciencias de la Computación, Escuela Politécnica Nacional, Quito, Ecuador

\vspace{0.5cm}
\hrule
\vspace{0.5cm}

\section*{Ethical Approval and Research Compliance}

\subsection*{1. Human Participants and Patient Data}
This study \textbf{did not involve human participants, clinical trials, medical interventions, or identifiable human tissues/specimens}. Consequently, institutional ethical review and formal informed consent procedures were \textbf{not required (exempt)} in accordance with the Declaration of Helsinki and international guidelines for computational biology research. The non-antimicrobial background sequences used for negative benchmarking were retrieved programmatically from the publicly accessible, anonymized, and reviewed UniProtKB/Swiss-Prot database (canonical proteome of \textit{Homo sapiens}).

\subsection*{2. Animal Subjects and Welfare}
This research \textbf{did not involve experimental work with live vertebrate or invertebrate animal subjects}. No animals were bred, handled, housed, or sacrificed for the purpose of this study. All primary peptide sequence records analyzed herein were curated programmatically from open-access scientific repositories (UniProtKB/Swiss-Prot release 2024\_02). Formal approval from an Institutional Animal Care and Use Committee (IACUC) was therefore not applicable.

\subsection*{3. Biosecurity and Dual-Use Research of Concern (DURC)}
The computational methods and algorithms introduced in this work (\texttt{AMP-LeakGuard}) are designed strictly for \textbf{benchmarking data leakage, evaluating negative sampling bias, and enhancing the statistical rigor of predictive models in bioinformatics}. 
\begin{itemize}
    \item The study does \textbf{not} provide protocols or instructions for the genetic engineering, de novo synthesis, or weaponization of harmful biological agents, Select Agents, or high-consequence pathogens.
    \item During dataset curation, known toxins (\texttt{KW-0800}) and harmful venoms were explicitly filtered out and excluded from the experimental design.
    \item The primary objective of this research is defensive and therapeutic: improving the computational discovery of narrow-spectrum antimicrobial agents against multi-drug resistant bacterial infections (AMR).
\end{itemize}

\subsection*{4. Data Integrity and Algorithmic Fairness}
All datasets, data processing scripts, machine learning training procedures, and statistical evaluation pipelines are fully reproducible and released under the permissive open-source MIT License (\url{https://github.com/MrOnie/AMP-LeakGuard}). No private, proprietary, or restricted-access databases were used in this investigation.

\vspace{1.0cm}

\noindent\textbf{Ronie Martínez}\\
Corresponding Author\\
Departamento de Informática y Ciencias de la Computación\\
Escuela Politécnica Nacional, Quito, Ecuador\\
Email: \href{mailto:ronie.martinez@epn.edu.ec}{ronie.martinez@epn.edu.ec}

\end{document}
